%%%%%
%
%% Dans ce fichier, les �ventuels accents sont cod�s
%% en ISO-8859-15 (aka latin9)
%% La ligne suivante donne trois exemples 
%% e accent aigu : �, e accent grave : �, e-dans-l'o : �
\section{Vitesse conventionnelle a�ronautique}

\itemstyle{Op\'erateur~:} \operfmt{aero\_vit\_conv}
\index{aero\_vit\_conv}

%%%%%OBJET%%%%
\itemstyle{Objet~:}
Calcul d'une vitesse conventionnelle a�ronautique � partir
d'une pression dynamique. Cette formule ne donne une
vraie vitesse que lors de mesures dans les 
conditions standard au niveau de la mer.

En appliquant cette formule � une mesure directe
d'un pitot d'un a�ronef, on obtient une \emph{vitesse (conventionnelle)
indiqu�e}, ou IAS. En appliquant cette formule � une mesure 
de pitot corrig�e des erreurs instrumentales connues (en particulier
les effets li�s � sa position : erreur 
"de statique"), on obtient une \emph{vitesse corrig�e}, ou CAS.


%%%%%ARGUMENTS%%%%
\bigskip
%% Dans ce fichier, les �ventuels accents sont cod�s
%% en ISO-8859-15 (aka latin9)
%% La ligne suivante donne trois exemples 
%% e accent aigu : �, e accent grave : �, e-dans-l'o : �
%%%%%ENTREES%%%%
\begin{tabarg}
\multicolumn{3}{l}{\titretab {Entr\'ees}}\\
\hline
$\Delta P$ & Vecteur, \texttt{float} & Pression dynamique : diff�rence entre la 
                  pression d'arr�t et la pression statique (hPa)\\
unit$_{out}$ & Coeff., \texttt{String} & Unit�e souhait�e pour le 
                  r�sultat (kt ou m/s)\\
\end{tabarg}

\medskip
%%%%%SORTIES%%%%
\begin{tabarg}
\multicolumn{3}{l}{\titretab {Sorties}}\\
\hline
$V_c$ & Vecteur, \texttt{float} & Vitesse conventionnelle 
                correspondant � $\Delta P$ (unit$_{out}$)\\
\end{tabarg}


%%%%%FORMULE%%%%
\itemstyle{Formule}

%% Formule de l'operateur : aero_vit_conv
%% Dans ce fichier, les �ventuels accents sont cod�s
%% en ISO-8859-15 (aka latin9)
%% La ligne suivante donne trois exemples 
%% e accent aigu : �, e accent grave : �, e-dans-l'o : �
Si unit$_{out}$ = "kt", alors $V_{s0} = 661.479$. Sinon, si
unit$_{out}$ = "m/s", alors  $V_{s0} = 340.294 $
\[
V_c = V_{s0} \sqrt{5 \left(\frac{\Delta P}{1013.25}\right)^{\frac{2}{7}}}
\]



%%%%%AUTEUR%%%%
\itemstyle{Auteur~:}
CNRM/GMEI/TRAMM

%%%%%REFERENCE%%%%
\itemstyle{R\'ef\'erences~:}

Norme AFNOR  NF X 02-115 : \emph{Symbole et vocabulaire de la m�canique
du vol} \cite{AFNOR-74}, paragraphe 5.6.1 page 45.

Cette norme fixe en particulier les valeurs num�riques de $P_0$ 
(1013.25 hPa), de la vitesse du son $a_n$ (340.294 m.s$^{-1}$),
et de $\gamma$ (1.400, ce qui induit $\frac{2}{\gamma-1} = 5$
et $\frac{\gamma-1}{\gamma} = 2/7$).
